\chapter{Fourier Series}

A Fourier series writes a periodic function as a sum of sines and cosines. We need it to fit initial conditions in PDEs. The eigenfunctions of Chapter 6 are exactly these sines and cosines.

\section{Definition and formulas}

\begin{defbox}{Fourier series of period $2L$}
Let $f$ be periodic with period $2L$. Its Fourier series is
\[ f(x) = \frac{a_0}{2} + \sum_{n=1}^{\infty}\Big(a_n\cos\frac{n\pi x}{L} + b_n\sin\frac{n\pi x}{L}\Big), \]
with coefficients
\[ a_n = \frac1L\int_{-L}^{L}f(x)\cos\frac{n\pi x}{L}\,\dd x \ (n\ge0), \qquad b_n = \frac1L\int_{-L}^{L}f(x)\sin\frac{n\pi x}{L}\,\dd x \ (n\ge1). \]
\end{defbox}

These formulas come from orthogonality (Chapter 6). Multiply both sides by $\cos(m\pi x/L)$ or $\sin(m\pi x/L)$, integrate over one period, and all terms but one vanish.

\begin{derivbox}{Euler's formulas for the coefficients}
Assume $f(x) = \frac{a_0}{2} + \sum\big(a_n\cos\frac{n\pi x}{L} + b_n\sin\frac{n\pi x}{L}\big)$ and that we may integrate term by term. We need these integrals over one period, for integers $m,n\ge1$:
\[ \int_{-L}^{L}\cos\frac{n\pi x}{L}\,\dd x = 0,\qquad \int_{-L}^{L}\sin\frac{n\pi x}{L}\,\dd x = 0, \]
\[ \int_{-L}^{L}\cos\frac{m\pi x}{L}\cos\frac{n\pi x}{L}\,\dd x = \begin{cases}L & m=n\\0 & m\ne n\end{cases},\qquad \int_{-L}^{L}\sin\frac{m\pi x}{L}\sin\frac{n\pi x}{L}\,\dd x = \begin{cases}L & m=n\\0 & m\ne n\end{cases}, \]
\[ \int_{-L}^{L}\sin\frac{m\pi x}{L}\cos\frac{n\pi x}{L}\,\dd x = 0\ \text{ for all } m,n. \]
They follow from $\cos A\cos B = \frac12[\cos(A-B)+\cos(A+B)]$, $\sin A\sin B = \frac12[\cos(A-B)-\cos(A+B)]$ and $\sin A\cos B = \frac12[\sin(A+B)+\sin(A-B)]$, exactly as in Chapter 6.

\textbf{Find $a_0$.} Integrate the series over $(-L,L)$. Every term vanishes except the constant: $\int_{-L}^{L}f\,\dd x = \frac{a_0}{2}\cdot2L$, so $a_0 = \frac1L\int_{-L}^{L}f\,\dd x$.

\textbf{Find $a_m$.} Multiply the series by $\cos\frac{m\pi x}{L}$ and integrate over $(-L,L)$. Only the term $a_m\cos\frac{m\pi x}{L}$ survives, and it gives $a_m\cdot L$. So $a_m = \frac1L\int_{-L}^{L}f(x)\cos\frac{m\pi x}{L}\,\dd x$.

\textbf{Find $b_m$.} Multiply by $\sin\frac{m\pi x}{L}$ and integrate: $b_m = \frac1L\int_{-L}^{L}f(x)\sin\frac{m\pi x}{L}\,\dd x$.

The formula for $a_n$ with $n=0$ would give twice the constant term. This is why the constant term is written as $a_0/2$.
\end{derivbox}

\begin{defbox}{Convergence (Dirichlet conditions, stated)}
If $f$ is piecewise smooth on one period (finite number of jumps and finite number of maxima and minima), then the series converges. At every point where $f$ is continuous the sum equals $f(x)$. At a jump, the sum equals the \textbf{average of the left and right limits}.
\end{defbox}

\begin{defbox}{Even and odd functions}
\begin{itemize}
  \item If $f$ is \textbf{even} ($f(-x)=f(x)$): all $b_n = 0$, and $a_n = \dfrac2L\displaystyle\int_0^L f(x)\cos\dfrac{n\pi x}{L}\,\dd x$. We get a \textbf{cosine series}.
  \item If $f$ is \textbf{odd} ($f(-x)=-f(x)$): all $a_n=0$, and $b_n = \dfrac2L\displaystyle\int_0^L f(x)\sin\dfrac{n\pi x}{L}\,\dd x$. We get a \textbf{sine series}.
\end{itemize}
Check for even or odd first. It saves half the work.
\end{defbox}

\begin{derivbox}{Even and odd functions, and half-range series}
\textbf{Odd $f$.} Then $f(x)\cos\frac{n\pi x}{L}$ is (odd)$\times$(even) $=$ odd, and the integral of an odd function over $(-L,L)$ is zero. So $a_n = 0$. The product $f(x)\sin\frac{n\pi x}{L}$ is even, so $\int_{-L}^{L} = 2\int_0^L$, and $b_n = \frac2L\int_0^Lf\sin\frac{n\pi x}{L}\,\dd x$.

\textbf{Even $f$.} In the same way $b_n = 0$ and $a_n = \frac2L\int_0^Lf\cos\frac{n\pi x}{L}\,\dd x$.

\textbf{Half-range series.} Given $f$ only on $(0,L)$, define $f_{\text{odd}}(-x) = -f(x)$ and $f_{\text{even}}(-x) = f(x)$ on $(-L,0)$. Each extension is a function on $(-L,L)$ with its own Fourier series (a sine series for the odd one and a cosine series for the even one). On $(0,L)$ both series equal the original $f$.
\end{derivbox}

\section{Worked examples}

\begin{examplebox}{Square wave}
$f(x) = -1$ on $(-\pi, 0)$ and $f(x)=1$ on $(0,\pi)$, period $2\pi$ (so $L=\pi$).

$f$ is odd, so $a_n = 0$ and
\[ b_n = \frac2\pi\int_0^\pi \sin nx\,\dd x = \frac{2}{\pi}\cdot\frac{1 - \cos n\pi}{n} = \begin{cases}\dfrac{4}{n\pi} & n \text{ odd},\\[4pt] 0 & n\text{ even}.\end{cases} \]
\[ f(x) = \frac4\pi\Big(\sin x + \frac{\sin 3x}{3} + \frac{\sin 5x}{5} + \cdots\Big). \]
At the jump $x=0$ the series gives $0$, the average of $-1$ and $1$.
\end{examplebox}

\begin{center}
\begin{tikzpicture}
\begin{axis}[width=0.92\textwidth,height=5.4cm,trig format plots=rad,xmin=-3.2,xmax=3.2,ymin=-1.5,ymax=1.5,
  xlabel={$x$},grid=both,grid style={gray!20},legend pos=south east,legend style={font=\small},
  xtick={-3.14159,0,3.14159},xticklabels={$-\pi$,$0$,$\pi$}]
\addplot[black!60,dashed,thick] coordinates {(-3.14159,-1) (-0.0001,-1) (0.0001,1) (3.14159,1)};
\addlegendentry{$f$}
\addplot[accent,thick,domain=-3.14159:3.14159,samples=200]{4/pi*(sin(x)+sin(3*x)/3)};
\addlegendentry{2 terms}
\addplot[notec,thick,domain=-3.14159:3.14159,samples=400]{4/pi*(sin(x)+sin(3*x)/3+sin(5*x)/5+sin(7*x)/7+sin(9*x)/9+sin(11*x)/11)};
\addlegendentry{6 terms}
\end{axis}
\end{tikzpicture}

{\small Figure: partial sums of the square-wave series. More terms follow the wave better, but there is always an overshoot of about 9\% near the jump (Gibbs phenomenon).}
\end{center}

\begin{examplebox}{The function $f(x)=x$ on $(-\pi,\pi)$}
$f$ is odd. $b_n = \dfrac2\pi\displaystyle\int_0^\pi x\sin nx\,\dd x$. Integrate by parts:
$\displaystyle\int_0^\pi x\sin nx\,\dd x = \Big[-\frac{x\cos nx}{n} + \frac{\sin nx}{n^2}\Big]_0^\pi = -\frac{\pi\cos n\pi}{n}$. So
\[ b_n = -\frac{2\cos n\pi}{n} = \frac{2(-1)^{n+1}}{n}, \qquad x = 2\Big(\sin x - \frac{\sin 2x}{2} + \frac{\sin 3x}{3} - \cdots\Big). \]
\end{examplebox}

\begin{examplebox}{The function $f(x)=x^2$ on $(-\pi,\pi)$}
$f$ is even. $a_0 = \dfrac2\pi\displaystyle\int_0^\pi x^2\,\dd x = \dfrac{2\pi^2}{3}$. For $n\ge1$,
\[ a_n = \frac2\pi\int_0^\pi x^2\cos nx\,\dd x = \frac2\pi\Big[\frac{x^2\sin nx}{n} + \frac{2x\cos nx}{n^2} - \frac{2\sin nx}{n^3}\Big]_0^\pi = \frac{4(-1)^n}{n^2}. \]
\[ x^2 = \frac{\pi^2}{3} + 4\sum_{n=1}^{\infty}\frac{(-1)^n}{n^2}\cos nx. \]
\textbf{A bonus.} Put $x = \pi$: $\pi^2 = \dfrac{\pi^2}{3} + 4\sum\dfrac1{n^2}$, so $\displaystyle\sum_{n=1}^\infty\frac{1}{n^2} = \frac{\pi^2}{6}$.
\end{examplebox}

\section{Half-range expansions}

Suppose $f$ is given only on $(0,L)$ (a rod or a string). We can extend it to $(-L,0)$ in two ways, and each choice gives a different series for the same $f$ on $(0,L)$.

\begin{methodbox}{Half-range series}
\begin{itemize}
  \item \textbf{Sine series} (odd extension): $f(x) = \displaystyle\sum_{n=1}^\infty b_n\sin\frac{n\pi x}{L}$, \quad $b_n = \dfrac2L\displaystyle\int_0^L f(x)\sin\frac{n\pi x}{L}\,\dd x$.
  \item \textbf{Cosine series} (even extension): $f(x) = \dfrac{a_0}{2} + \displaystyle\sum_{n=1}^\infty a_n\cos\frac{n\pi x}{L}$, \quad $a_n = \dfrac2L\displaystyle\int_0^L f(x)\cos\frac{n\pi x}{L}\,\dd x$.
\end{itemize}
Use the sine series when the PDE has zero values at the ends ($u=0$). Use the cosine series when the ends are insulated ($u_x=0$).
\end{methodbox}

\begin{examplebox}{Cosine series of $f(x) = x$ on $(0,\pi)$}
$a_0 = \dfrac2\pi\displaystyle\int_0^\pi x\,\dd x = \pi$. For $n\ge1$:
$a_n = \dfrac2\pi\Big[\dfrac{x\sin nx}{n} + \dfrac{\cos nx}{n^2}\Big]_0^\pi = \dfrac{2\big((-1)^n - 1\big)}{\pi n^2}$, which is $-\dfrac{4}{\pi n^2}$ for odd $n$ and $0$ for even $n$.
\[ x = \frac\pi2 - \frac4\pi\Big(\cos x + \frac{\cos 3x}{9} + \frac{\cos 5x}{25} + \cdots\Big),\qquad 0<x<\pi. \]
Check at $x=0$: $0 = \dfrac\pi2 - \dfrac4\pi\cdot\dfrac{\pi^2}{8} = 0$. \checkmark
\end{examplebox}

\section{Parseval's identity}

\begin{defbox}{Parseval's identity}
The energy of $f$ equals the sum of the energies of its harmonics:
\[ \frac1L\int_{-L}^{L}\big[f(x)\big]^2\,\dd x = \frac{a_0^2}{2} + \sum_{n=1}^\infty\big(a_n^2 + b_n^2\big). \]
\end{defbox}

\begin{derivbox}{Parseval's identity}
Multiply the series $f = \frac{a_0}{2} + \sum(a_n\cos + b_n\sin)$ by $f$ itself and integrate over $(-L,L)$. By Euler's formulas, $\int f\,\dd x = La_0$, $\int f\cos\frac{n\pi x}{L}\,\dd x = La_n$ and $\int f\sin\frac{n\pi x}{L}\,\dd x = Lb_n$. So
\[ \int_{-L}^{L}f^2\,\dd x = \frac{a_0}{2}\,(La_0) + \sum_{n=1}^\infty\big(a_n\cdot La_n + b_n\cdot Lb_n\big). \]
Divide by $L$ to get $\dfrac1L\displaystyle\int_{-L}^{L}f^2\,\dd x = \dfrac{a_0^2}{2} + \sum\big(a_n^2 + b_n^2\big)$.
\end{derivbox}

For the square wave: the left side is $\frac1\pi\int_{-\pi}^{\pi}1\,\dd x = 2$. The right side is $\sum_{n\text{ odd}}\frac{16}{\pi^2n^2} = \frac{16}{\pi^2}\cdot\frac{\pi^2}{8} = 2$. \checkmark

\begin{practicebox}
\begin{enumerate}
  \item Find the Fourier series of $f(x)=|x|$ on $(-\pi,\pi)$.
  \item Find the Fourier series of $f(x)=x$ on $(-1,1)$ (period 2, so $L=1$).
  \item Find the half-range sine series of $f(x) = 1$ on $(0,L)$.
  \item Use the series of $x^2$ at $x=0$ to find $\displaystyle\sum_{n=1}^\infty\frac{(-1)^{n+1}}{n^2}$.
  \item To what value does the square-wave series converge at $x=0$ and at $x=\pi/2$?
\end{enumerate}
\end{practicebox}

\begin{answerbox}
(1) Even. $a_0 = \pi$, $a_n = \dfrac{2((-1)^n - 1)}{\pi n^2}$: $|x| = \dfrac\pi2 - \dfrac4\pi\Big(\cos x + \dfrac{\cos 3x}{9} + \cdots\Big)$. \quad (2) $f = \dfrac2\pi\displaystyle\sum_{n=1}^\infty\frac{(-1)^{n+1}}{n}\sin n\pi x$. \quad (3) $b_n = \dfrac{2}{L}\displaystyle\int_0^L\sin\dfrac{n\pi x}{L}\,\dd x = \dfrac{2(1-(-1)^n)}{n\pi}$, so $f = \dfrac4\pi\displaystyle\sum_{n\text{ odd}}\dfrac1n\sin\dfrac{n\pi x}{L}$. \quad (4) $0 = \frac{\pi^2}{3} + 4\sum\frac{(-1)^n}{n^2}$, so $\sum\frac{(-1)^n}{n^2} = -\frac{\pi^2}{12}$ and the answer is $\dfrac{\pi^2}{12}$. \quad (5) At $x = 0$ (a jump) it converges to $0$. At $x=\pi/2$ (continuous) it converges to $1$.
\end{answerbox}
